\section{Integral Convergence}
\label{sec:convergenza_integrale}

Motivated by the way the norm of a vector in $\RR^n$ is defined, it is natural
to give the following definition of the \emph{Euclidean norm} for a function
$f\colon (a,b)\to \RR$:
\[
  \Abs{f}_2 = \sqrt{\int_a^b \abs{f(x)}^2\, dx}.
\]
This definition makes sense, as an improper integral, if the function $f$ is
locally Riemann-integrable on the interval $(a,b)$
(definition~\ref{def:localmente_riemann}). In such a case, the integral exists,
but it might take the value $+\infty$. We then define
\[
\H(a,b) =
\ENCLOSE{\text{$f$ locally Riemann-integrable on $(a,b)$}\colon \Abs{ f}_2< +\infty}.
\]

We will now attempt to demonstrate that the norm we have introduced is indeed a
norm (in the sense of definition~\ref{def:norma}) that renders $\H(a,b)$ an
infinite-dimensional Euclidean vector space.

First of all, it is clear that for any $f$, we have $\Abs{ f}_2\ge 0$. Moreover,
it is trivial to observe that the norm is homogeneous: if $t\in \RR$, we have:
\[
  \Abs{t\cdot f}_2= \abs{t}\cdot \Abs{f}_2.
\]
Thus, if $f\in \H(a,b)$, then $t f\in \H(a,b)$ for every $t \in \RR$. If $f,g\in
\H(a,b)$, by the parallelogram law~\eqref{eq:parallelogramma} valid in $\RR$, we
have
\[
  (f(x) + g(x))^2 \le 2f^2(x) + 2 g^2(x)
\]
and thus if $f,g \in \H(a,b)$, then $f+g\in \H(a,b)$. Therefore, $\H(a,b)$ is a
vector space.

Note that if $f$ and $g$ are locally Riemann-integrable, then $f+g$ is also
locally Riemann-integrable, and thanks to theorem~\ref{th:composta_integrabile},
so are $f^2$, $g^2$, and $(f+g)^2$. Consequently, $f\cdot g =
\frac{(f+g)^2-f^2-g^2}{2}$ is Riemann-integrable. For $f,g\in \H(a,b)$, we can
then define the inner product:
\begin{equation}\label{eq:854392}
  \langle f, g\rangle = \int_a^b f(x) \cdot g(x)\, dx.
\end{equation}
Young's inequality~\eqref{eq:Young} valid in $\RR$ tells us that
\[
  \abs{f(x) \cdot g(x)} \le \frac{f^2(x)+g^2(x)}{2}
\]
and thus ensures that the integral~\eqref{eq:854392} is absolutely convergent.
Therefore, if $f,g\in \H(a,b)$, the inner product $\langle f,g\rangle$ is
well-defined and is a finite number. At this point, it is easy to verify that
$\langle f,g\rangle$ is a symmetric non-negative bilinear form. Thus, it
satisfies all the formal properties given in
definition~\ref{def:prodotto_scalare} of an inner product, except that it is not
guaranteed that $\langle f,f\rangle = 0$ implies $f=0$. Indeed, this property is
false because if the function $f^2(x)$ has a zero integral, it is not
necessarily identically zero: the integral, in fact, does not change if the
function is modified at a single point.

To address this issue, we must identify two functions $f,g\in \H(a,b)$ if
$\Abs{f-g}_2=0$:
\[
  f \sim g  \iff \Abs{f-g}_2=0.
\]
We call $H(a,b)$ the quotient:
\[
  H(a,b) = \H(a,b)/\sim
\]
that is, the vector space of functions in $\H(a,b)$ where two functions are
identified if the norm of their difference is zero.

The quotient space $H(a,b)$ is finally an Euclidean space. It indeed retains the
vector space structure (since the set of $f\in \H(a,b)$ with $\Abs{f}_2=0$ is a
vector space), and the inner product $\langle f,g\rangle$ (and consequently the
norm $\Abs{f}_2$) is well-defined on $H(a,b)$, and in addition to the properties
we have already demonstrated, we can now assert that $\langle f,f\rangle =
\Abs{f}_2=0$ if and only if $f=0$.

\begin{remark}
The notation $H(a,b)$ that we choose to use for defining this space is not
standard. The letter $H$ refers to Hilbert, as the study of this particular
space led to the abstraction of the concept of a Hilbert space. What we will see
in the following is that this space is not complete and thus is not actually a
Hilbert space. In fact, the use of improper integrals is an attempt to extend
the Riemann integral to make this space complete. Only with the definition of
the Lebesgue integral (1875--1941) was it possible to extend the Riemann
integral to a broader class of functions to complete the space. The analogue of
the space $H(a,b)$ defined via the Lebesgue integral is usually called $L^2$.
The exponent $2$ is due to the fact that it would be possible to define $L^p$
spaces in a manner analogous to what was done in example~\ref{ex:norma_p}. It
would still hold that only for $p=2$ are these normed spaces Euclidean, that is,
only for $p=2$ is the norm induced by an inner product. All $L^p$ spaces are
complete and thus are Banach spaces. But $L^2$ is also a Hilbert space, as it is
equipped with an inner product. Indeed, we could say that $L^2$ is the
quintessential Hilbert space and is also denoted by the name $H^0$ (in this
case, the exponent denotes the differentiability of the functions, as in $C^0$).
\end{remark}

The elements of $H(a,b)$ are not functions, but equivalence classes of
functions. In the following, however, we will treat $f\in H(a,b)$ as if it were
$f\in \H(a,b)$, ensuring that our statements remain valid if, instead of a
function $f$, an equivalent function is taken. In particular, it will make sense
to consider the integrals of these functions because the value of the integral
does not depend on the chosen representative, but it will not make sense to
consider the value that the functions take at a single point, as this depends on
the representative.

\subsection{Fourier Series}
\index{Fourier!series}%
\index{series!Fourier}%

The space $H(a,b)$ is a real vector space on which we have managed to define an
inner product and thus a norm. We would like to find an orthonormal basis on
$H(a,b)$ with respect to which it is possible to write the coordinates of the
vectors (i.e., the functions) of $H(a,b)$. Since $H(a,b)$ does not have finite
dimension, we cannot hope to find a basis with a finite number of elements. We
therefore introduce the following.

\begin{definition}[Hilbertian basis]
\index{Hilbert!basis}%
Let $V$ be an Euclidean space. Let $e_k\in V$ with $k\in \NN$. We will say that
$e_k$ is an \emph{orthonormal system}%
\mymargin{orthonormal system}\index{orthonormal system} if $\langle e_k, e_j\rangle = 0$ when $k\neq j$ and $\langle e_k, e_k\rangle = 1$ for every $k\in \NN$. We will say that $e_k$ is a \emph{Hilbertian basis}%
\mymargin{Hilbertian basis}\index{basis!Hilbertian} if it is an orthonormal system and moreover for every $v\in V$, there exist $c_k\in \RR$ with $k\in\NN$ such that
\begin{equation}\label{eq:49346}
  v = \sum_{k=0}^{+\infty} c_k e_k.
\end{equation}
\end{definition}

\begin{theorem}[Bessel's inequality]
\label{th:bessel}%
Let $e_0, e_1, e_2, \dots$ be an orthonormal system in an Euclidean space $V$.
For every $v\in V$ and for every $k\in \NN$, let $c_k = \langle v,e_k\rangle$.
Then the \emph{Bessel's inequality}%
\mymargin{Bessel's inequality}\index{inequality!Bessel's} holds:
\index{Bessel!inequality of}
\begin{equation}\label{eq:disuguaglianza_Bessel}
\sum_{k=0}^{+\infty} c_k^2 \le \abs{v}^2.
\end{equation}

Moreover, the equality
\begin{equation}\label{eq:487562}
\sum_{k=0}^{+\infty} c_k^2 = \abs{v}^2
\end{equation}
holds if and only if \eqref{eq:49346} holds.

We can therefore assert that the orthonormal system $e_0, e_1, e_2, \dots$ is a
Hilbertian basis if and only if for every $v\in V$, the
equality~\eqref{eq:487562} holds.
\end{theorem}
%
\begin{proof}
Given $v\in V$, let $c_k = \langle v,e_k\rangle$ and
\[
   v_N = \sum_{k=0}^N c_k e_k, \qquad w_N = v - v_N.
\]
Observe that for every $n\in \NN$, we have $\langle v_N,w_N\rangle=0$ ($w_N$ is
orthogonal to $v_N$) because if $n\le N$:
\[
  \langle w_N , e_n \rangle
  = \langle v, e_n\rangle - \sum_{k=0}^N c_k \langle e_k, e_n\rangle
  = \langle v,e_n\rangle - c_n = 0
\]
and therefore
\[
 \langle w_N, v_N \rangle = \sum_{k=0}^N c_k \langle w_N, e_k\rangle = 0.
\]
Thus, we can conclude, thanks to the Pythagorean theorem~\eqref{eq:Pitagora}
\begin{align*}
  \abs{v}^2
  &= \abs{v_N + w_N}^2
  = \abs{v_N}^2 + \abs{w_N}^2 \\
  &\ge \abs{v_N}^2 = \abs{\sum_{k=0}^N c_k e_k}^2
  = \sum_{k=0}^N c_k^2.
\end{align*}
Taking the limit as $N\to +\infty$, we obtain Bessel's inequality.

We have seen that $\abs{w_N}^2 = \abs{v}^2 - \abs{v_N}^2$, thus the condition
$\abs{w_N}^2\to 0$ is equivalent to Bessel's equality: $\abs{v_N}^2 \to
\abs{v}^2$. The system is Hilbertian, by definition, if and only if $v_N \to v$,
i.e., if $w_N\to 0$.
\end{proof}

Our goal is now to produce a Hilbertian basis for $H(0,2\pi)$. A basis for
$H(a,b)$ can be obtained consequently by appropriately translating and rescaling
the functions of the basis of $H(0,2\pi)$. The previous theorem tells us that if
$H(a,b)$ admits a Hilbertian basis $e_k$, then every function $u\in H(a,b)$ can
be simply represented by a numerical sequence $x_k = (u, e_k)$ such that
\[
    \Abs{u}_2 = \sqrt{\sum_k x_k^2}.
\]
\mynote{%
If $u$ represents the intensity of the acoustic signal emitted by a musical
instrument playing a certain note, we can imagine that $u$ is a periodic
function with frequency $\nu$ equal to the frequency of the note emitted. The
functions $e_k$ of the Hilbertian basis are sinusoidal functions, each with
frequency $\nu k$. The coefficients $x_k$ in the decomposition $u = \sum x_k
u_k$ represent the intensities of each of these fundamental frequencies and
determine the \emph{timbre} of the musical instrument.
}

Consider the following trigonometric functions%
\mynote{%
From an algebraic point of view, it would be much simpler to consider functions
with complex values:
\[
  e_k(x) = \frac{e^{ikx}}{\sqrt{2\pi}}
\]
varying $k\in \ZZ$. However, we prefer to remain within the realm of real
functions to avoid introducing the Hermitian structure of complex vector
spaces.}%
:
\begin{equation}\label{eq:54741346}
\begin{aligned}
  e_0 (x) &= \frac{1}{\sqrt{2\pi}} \\
  e_{2k+1}(x) &= \frac{\sin(kx)}{\sqrt{\pi}} \qquad k=0,1,\dots\\
  e_{2k}(x) &= \frac{\cos(kx)}{\sqrt{\pi}} \qquad k=1,2,\dots
\end{aligned}
\end{equation}
Obviously
\[
  \Abs{e_0}_2=\sqrt{\int_{0}^{2\pi} \enclose{\frac{1}{\sqrt{2\pi}}}^2\, dx} = 1
\]
but it is also easy to verify that
\[
  \int_{0}^{2\pi} \cos^2 (kx)\, dx
  = \frac{1}{k}\int_{0}^{2k\pi} \cos^2 y\, dy
  = \int_{0}^{2\pi} \cos^2 y\, dy = \pi
\]
from which we obtain $\Abs{e_{2k}}_2 = 1$. Since the sine function is the
translation of the cosine function, we similarly obtain that
$\Abs{e_{2k+1}}_2=1$. Thus, the vectors $e_0,e_1,\dots$ are all of unit modulus.
Using Euler's formula, one can find the \emph{Werner's formula}%
\mymargin{Werner's formula}\index{Werner's formula}:
\index{Werner!formula of}%
\index{formula!Werner's}%
\begin{align*}
 \sin(mx)\cos(nx)
 &= \frac{e^{imx}-e^{-imx}}{2i}\cdot \frac{e^{inx}+e^{-inx}}{2}\\
 &= \frac{e^{i(m+n)x}}{4i} - \frac{e^{-i(m+n)x}}{4i} 
  + \frac{e^{i(m-n)x}}{4i} - \frac{e^{-i(m-n)x}}{4i}\\
 &= \frac{\sin\enclose{(m+n)x}}{2} + \frac{\sin\enclose{(m-n)x}}{2}
\end{align*}
from which, if $m\neq n$:
\begin{align*}
 \int_{0}^{2\pi} \sin(mx)\cos(nx)\, dx &=
 -\frac 1 {2(m+n)} \Enclose{\cos\enclose{(m+n)x}}_{0}^{2\pi} \\
  &\quad - \frac{1}{2(m-n)}\Enclose{\cos\enclose{(m-n)x}}_{0}^{2\pi} = 0
\end{align*}
and if $m=n$, the result is still the same. This means that $\langle
e_{2n},e_{2m+1}\rangle = 0$. A similar argument can be made for the functions
$\sin(mx)\sin(nx)$ and $\cos(mx)\cos(nx)$, finding, in those cases as well, that
such products have zero integral over the interval $[0,2\pi]$. For our memory,
the Werner's formulas found in these last cases are:
\mynargin{Werner's formulas}%
\index{Werner's formulas}%
\index{Werner!formula of}%
\index{formula!Werner's}%
\begin{align*}
  \cos(mx) \cos(nx) &=  \frac{\cos((m+n)x)}{2} + \frac{\cos((m-n)x)}{2}\\
  \sin(mx) \sin(nx) &=  \frac{\cos((m-n)x)}{2} - \frac{\cos((m+n)x)}{2}.
\end{align*}

It follows that the functions $e_0,e_1,e_2,\dots$ defined by~\eqref{eq:54741346}
form an orthonormal system in $H(0,2\pi)$ since we have:
\[
  \langle e_n, e_m \rangle =
  \begin{cases} 1 &\text{if $m=n$}\\
  0 & \text{otherwise}.
  \end{cases}
\]
Consequently, the functions $e_0,e_1,e_2, \dots$ are linearly independent in
$H(0,2\pi)$ because if a linear combination of them is zero:
\[
  \sum_{k=0}^n \lambda_k e_k = 0
\]
then, taking the inner product with $e_j$, we obtain:
\[
  0 = \langle \sum_{k=0}^n \lambda_k e_k , e_j\rangle
  = \sum_{k=0}^n \lambda_k \langle e_k, e_j\rangle
  = \lambda_j
\]
and thus all the coefficients $\lambda_j$ are zero. We can, in particular,
deduce that $H(0,2\pi)$ has infinite dimension.

\begin{definition}[trigonometric polynomials]
A function $f$ is said to be a \emph{trigonometric polynomial}%
\mynargin{trigonometric polynomial}\index{trigonometric polynomial} if there exists a polynomial in two variables $P(X,Y)$ such that
\[
  f(x) = P(\cos x, \sin x).
\]
\end{definition}

We can verify that $f$ is a trigonometric polynomial if and only if it is
possible to write $f$ as a finite linear combination of the functions $e_k$ just
introduced, that is, if there exist $c_0, c_1, \dots, c_n\in \RR$ such that
\[
  f = \sum_{k=0} c_k e_k.
\]
It suffices to reduce the trigonometric functions to the complex exponential,
using Euler's formula:
\begin{align*}
   \cos(nx) + i \sin(nx)
   &= e^{inx} = (e^{ix})^n
   = \enclose{\cos x + i \sin x}^n\\
   &= \sum_{k=0}^n {n \choose k} i^k \enclose{\sin x}^k \enclose{\cos x}^{n-k}
\end{align*}
from which, taking the real part and the imaginary part, one can express
$\cos(nx)$ and $\sin(nx)$ as polynomials in $\sin x$ and $\cos x$. More
precisely, for the cosine, only the terms with even $k$ are obtained:
\begin{align*}
\cos(nx) &= \sum_{k=0}^{\floor{\frac n 2}} {n \choose 2k} (-1)^k \enclose{\sin x}^{2k} \enclose{\cos x}^{n-2k}\\
&= \sum_{k=0}^{\floor{\frac n 2}}{n \choose 2k} \enclose{\cos^2 x -1}^k
\enclose{\cos x}^{n-2k}
\end{align*}
that is,
\[
  \cos(nx) = T_n(\cos x)
\]
where
\[
    T_n(x) = \sum_{k=0}^{\floor{\frac n 2}}{n \choose 2k} \enclose{x^2 -1}^k
  x^{n-2k}
\]
is called the \emph{Chebyshev polynomial of the first kind}.
\index{polynomial!Chebyshev}%
\index{Chebyshev!polynomial of}%
This means that each of the functions $e_k$ we introduced in~\eqref{eq:54741346}
is a trigonometric polynomial, and therefore any finite combination of such
functions is still a trigonometric polynomial.

Conversely, to verify that every trigonometric polynomial can be represented as
a finite linear combination of the elements $e_k$, we can observe that if $f(x)
= P(\cos x,\sin x)$ is a trigonometric polynomial, then using Euler's formulas,
one can write $f(x) = Q(e^{ix}, e^{-ix})$ where $Q$ is a polynomial with complex
coefficients. But observing that
\[
  \enclose{e^{ix}}^n = e^{inx}
\]
one obtains that $f(x)$ can be written as a complex linear combination of the
functions $e^{ikx}$ and $e^{-ikx}$. Using Euler's formulas again, one deduces
that $f(x)$ can be written as a complex linear combination of the functions
$\cos(ikx)$ and $\sin(ikx)$:
\[
  f(x) = \sum_{k=0}^n \enclose{z_k \cos(kx) + w_k \sin(kx)}.
\]
But since we know that $f(x)$ is a real function, the sum of the imaginary parts
of this linear combination is zero, and thus we have obtained that $f(x)$ is a
real linear combination of the functions $e_n$.

\begin{definition}[Fourier series]
If $f\in \H(0,2\pi)$, the coefficients $c_k = \langle f, e_k\rangle$ are called
the \emph{Fourier coefficients}%
\mynargin{Fourier coefficients}\index{Fourier coefficients} of $f$, and the series
\[
  \sum_{k=0}^{+\infty} c_k e_k
\]
is called the \emph{Fourier series}%
\mynargin{Fourier series}\index{Fourier series} of $f$. The trigonometric polynomial
\[
  f_N = \sum_{k=0}^{N} c_k e_k
\]
is called the \emph{Fourier expansion}%
\mynargin{Fourier expansion}\index{Fourier expansion} of order $N$ for $f$.
\end{definition}

\begin{example}
Let us calculate the Fourier expansion of the function:
\[
  f(x) = 
  \begin{cases}
    1 & \text{if $x\in [0,\pi]$}\\
    -1 & \text{if $x\in (\pi,2\pi]$}.
  \end{cases}
\]
Regarding $c_0 = \langle f, e_0\rangle$, we have trivially: 
\[
c_0 = \int_0^{2\pi} f(x) \cdot \frac{1}{\sqrt{2\pi}}\, dx 
  = \frac{1}{\sqrt{2\pi}}\Enclose{\int_0^\pi 1 \, dx + \int_0^\pi (-1)\, dx}
  = 0.  
\]
For $c_{2k} = \langle f, e_{2k}\rangle$, we have 
\begin{align*}
  c_{2k}
  &= \int_0^{2\pi} f(x)\frac{\cos (kx)}{\sqrt{\pi}}\, dx \\
  &= \frac{1}{\sqrt \pi} \Enclose{\int_0^\pi \cos (kx)\, dx 
  - \int_{\pi}^{2\pi} \cos(kx)\, dx}
  = 0 
\end{align*}
for reasons of symmetry.
To calculate $c_{2k+1}=\langle f, e_{2k+1}\rangle$, we find 
\begin{align*}
  c_{2k+1}
  &= \int_0^\pi \frac{\sin(kx)}{\sqrt \pi}\, dx 
    - \int_\pi^{2\pi} \frac{\sin(kx)}{\sqrt \pi}\, dx \\
  &= \frac{2}{\sqrt \pi}\int_0^\pi \sin(kx)\, dx
  = \frac{2}{k\sqrt \pi}\int_0^{k\pi}\sin t\, dt\\
  &= 
  \begin{cases}
    0 & \text{if $k$ is even,}\\
    \frac{4}{k\sqrt \pi} & \text{if $k$ is odd.}
  \end{cases}
\end{align*}
We then find that   
\[
 f_{2n+1}(x)
 = \sum_{k=0}^n \frac{4}{(2k+1)\sqrt \pi}\sin((2k+1)x).
\]
See figure~\ref{fig:fourier}.
\end{example}

\begin{figure}
  \begin{center}
  \includegraphics[height=5cm]{fourier.png}
  \end{center}
  \captionsetup{singlelinecheck=off}
  \caption{The Fourier expansion of the function
  that equals $1$ on $[0,\pi]$ and $-1$ on $[\pi,2\pi]$:
  $
  f_{61}(x)=
  \frac{4\sin(x)}{\pi} + \frac{4\sin(3x)}{3\pi} + \frac{4\sin(5x)}{5\pi} + \frac{4\sin(7x)}{7\pi} + \frac{4\sin(9x)}{9\pi} + \frac{4\sin(11x)}{11\pi} + \frac{4\sin(13x)}{13\pi} + \frac{4\sin(15x)}{15\pi} + \frac{4\sin(17x)}{17\pi} + \frac{4\sin(19x)}{19\pi} + \frac{4\sin(21x)}{21\pi} + \frac{4\sin(23x)}{23\pi} + \frac{4\sin(25x)}{25\pi} + \frac{4\sin(27x)}{27\pi} + \frac{4\sin(29x)}{29\pi} + \frac{4\sin(31x)}{31\pi}
  $.
  The code to calculate the coefficients and generate the graph
  is found on page~\pageref{code:Fourier}.
  It can be noted that at the jump point $x=\pi$, 
  the Fourier expansion has an oscillation greater than the original jump 
  (Gibbs phenomenon).
  \index{Gibbs!phenomenon of}
  \index{phenomenon of Gibbs}
  The increase in oscillation is about 9\% and is independent 
  of the number of terms used in the Fourier expansion.
%  This means that the Fourier series, while converging in norm $\Abs{\cdot}_2$,
%  does not converge uniformly.
  }
  \label{fig:fourier}
  \end{figure}
    
It is natural to ask whether what we have introduced is a \emph{Hilbertian
basis} because only in such a case do the Fourier expansions approximate the
function. To achieve this result, we must assume the validity of the following
theorem, which would now be too lengthy to prove.

\begin{theorem}[density of trigonometric polynomials]
For every continuous function $f\in C^0([0,2\pi])$ with $f(0)=f(2\pi)$, there
exists a sequence $g_k\in \H(0,2\pi)$ of trigonometric polynomials (i.e., finite
linear combinations of the system $\ENCLOSE{e_k}$) such that $g_k$ converges
uniformly to $f$.

In particular, the set of trigonometric polynomials is dense, with the uniform
norm, in the vector subspace
\[
  C^0_0([a,b]) 
  = \ENCLOSE{u\in C^0([a,b])\colon u(a)=u(b)=0}
  \subset C^0([a,b]).
\]
\end{theorem}

Observe that uniform convergence is stronger than convergence in $H(a,b)$
because we have, trivially:
\[
  \Abs{f}_2^2 = \int_a^b f(x)^2\, dx
  \le \int_a^b \Abs{f}_\infty^2\, dx
  = (b-a) \cdot \Abs{f}_\infty^2.
\]
Thus, the previous theorem guarantees that every continuous function in
$\H(a,b)$ can be approximated, with respect to the norm $\Abs{\cdot}_2$, by
trigonometric polynomials. The following theorem allows us to extend this
property to all functions in $\H(a,b)$.

\begin{theorem}[density of continuous functions]
Given any $f\in H(a, b)$, for every $\eps>0$, there exists a continuous function
$g\colon [a,b]\to \RR$ with $g(a)=g(b)=0$ such that $\Abs{f-g}_2<\eps$.

In other words: $C^0_0([a,b])=\ENCLOSE{u\in C^0\colon u(a)=u(b)=0}$ is a dense
subspace in $H(a,b)$, that is, a set whose closure (in the topology of $H(a,b)$)
is all of $H(a,b)$.
\end{theorem}
%
\begin{proof}
\emph{Step 1. Suppose $f$ is bounded on $(a,b)$.}
In this case, we can think of $f$ as being defined on $[a,b]$, and the integral
is a "proper" Riemann integral. Let $M=\sup \abs{f}([a,b])$. For the conditions
of integrability, we know that for every $\eps>0$, it is possible to find a
subdivision $a=x_0 < x_1 < \dots < x_n =b$ on which the integral of $f$ differs
from the upper and lower integrals by less than $\frac{\eps}{2M}$. We can take
as the function $g$ a linear interpolation, that is, a function such that
$g(x_k) = f(x_k)$ on the points of the subdivision and that is linear on each
subinterval $[x_k,x_{k+1}]$. The function $g$ thus defined is bounded, on each
subinterval, between the $\inf$ and the $\sup$ of $f$, and thus we will have:
\[
 \int_a^b \abs{f(x)-g(x)}\, dx < \frac{\eps}{M}.
\]
However, since $\Abs{f}_\infty\le M$, we also have $\Abs{g}_\infty\le M$, and
therefore $\Abs{f-g}_2\le 2M$. But then
\[
  \Abs{f-g}_2^2 = \int_a^b \abs{f-g}^2
  \le \int_a^b \Abs{f-g}_\infty \cdot \abs{f(x)-g(x)}\, dx
  \le 2M \frac{\eps}{2M} = \eps.
\]

\emph{Step 2. General case.}
If the function is absolutely integrable in the improper sense on $(a,b)$ but is
not bounded, by the definition of improper integral, for every $\eps>0$, we can
find an interval $[\alpha,\beta]\subset (a,b)$ on which $f^2$ (and thus $f$) is
bounded, and the integral of $f^2$ on $(a,\alpha]$ and on $[\beta,b)$ is less
than $\eps$. Let $M=\sup \abs{f}([\alpha,\beta])$. Applying the previous step on
the interval $[\alpha,\beta]$, we obtain a continuous function $g$ defined on
$[\alpha,\beta]$ such that $\abs{g(x)}\le M$ and $\int\abs{f-g}^2\le \eps$. The
idea is now to take a point $\alpha'\in(a,\alpha)$ and a point
$\beta'\in(\beta,b)$ and extend the function $g$ over the entire $[a,b]$ by
setting $g=0$ on the intervals $[a,\alpha']$ and $[\beta',b]$ and making a
linear interpolation on the intervals $[\alpha',\alpha]$ and $[\beta,\beta']$.

Using the inequality $(f-g)^2 \le (f-g)^2+(f+g)^2 = 2f^2+2g^2$ (otherwise, one
could use Cauchy-Schwarz twice) and choosing $\alpha'$ very close to $\alpha$ so
that $\alpha-\alpha'<\eps/M^2$, we obtain
\[
  \int_a^\alpha \abs{f-g}^2 
  = \int_a^{\alpha'} f^2 + \int_{\alpha'}^\alpha \abs{f-g}^2
  \le \eps + 2\eps + 2M^2\cdot(\alpha-\alpha')
  \le 5\eps.
\]
A similar inequality will hold on the interval $[\beta,b]$. Thus, we obtain in
conclusion
\[
\int_a^b \abs{f-g}^2
\le 5\eps + \eps + 5\eps  = 11 \eps
\]
and the theorem is proven with $11\eps$ instead of $\eps$.
\end{proof}

Thanks to the two previous theorems, we know that every $f\in \H(a,b)$ can be
approximated by trigonometric polynomials $g_N\to f$, that is, by functions of
the form:
\[
  g_N = \sum_{k=0}^{N} a_{k,N} e_k.
\]
Here, we are assuming that the trigonometric polynomial $g_N$ has an order not
exceeding $N$, but this can always be assumed by discarding or repeating the
terms of the sequence.

If we now consider the Fourier expansions:
\begin{equation}\label{eq:48948784}
  f_N = \sum_{k=0}^N c_k e_k, \qquad c_k = \langle f,e_k\rangle
\end{equation}
we remember that for every $k\le N$, we have
\[
   \langle f-f_N, e_k \rangle
   = \sum_{k=0}^N \enclose{\langle f,e_k\rangle - c_k} = 0
\]
and therefore
\[
  \langle f-f_N , g_N -f_N\rangle
  = \sum_{k=0}^N \langle f-f_N, (a_{k,N}-c_k) e_k \rangle
  = 0.
\]
Thus, applying the Pythagorean theorem~\ref{th:spazio_euclideo},
\[
  \Abs{f-f_N}_2^2 = \Abs{f-g_N}_2^2 - \Abs{f_N-g_N}_2^2
  \le \Abs{f-g_N}_2^2 \to 0.
\]

Thanks to theorem~\ref{th:bessel}, we have thus obtained the following.
\begin{theorem}[Fourier series]
If $f\in \H(0,2\pi)$ and $c_k$, $f_N$ are defined by~\eqref{eq:48948784}, then
$f_k\to f$ in norm $\Abs{\cdot}_2$. In particular, Bessel's equality holds
\begin{equation}\label{eq:uguaglianza_Bessel}
  \sum_{k=0}^{+\infty} c_k^2 = \Abs{f}_2^2.
\end{equation}
\end{theorem}


This means, therefore, that the orthonormal system we have introduced is a
Hilbertian basis and thus that every function $f\in \H(0,2\pi)$ is approximated
with respect to the norm $\Abs{\cdot }_2$ by the Fourier expansions.

Recall that if $\Abs{f_N-f}_2 \to 0$, it is not necessarily the case that
pointwise convergence $f_N(x)\to f(x)$ holds for every $x\in (0,2\pi)$ because
the limit is unique in $H(a,b)$ but not in $\H(a,b)$, and thus it is possible
(and indeed can happen) that the pointwise limit of the Fourier expansions
differs, at some points, from the function being approximated.

Unfortunately, the space $H(a,b)$ is not complete, as can be seen in the
following example.

\begin{example}[thickened rationals]
\label{ex:incompletezza_di_H}%
We demonstrate that the space $H(0,1)$ is not complete.

We know that the set $[0,1]\cap \QQ$ is countable, so there exists a sequence
$q_k$ that lists all the rational numbers in the interval $[0,1]$. Let $r_k =
\frac{1}{4\cdot 2^k}$ and consider the intervals $I_k = [q_k-r_k,q_k+r_k]$. Take
the sequence of functions $f_n\colon [0,1]\to \RR$ defined by
\[
  f_n(x) =
  \begin{cases}
  1 &\text{if } x\in\displaystyle\bigcup_{k=1}^{n} I_k\\
  0 & \text{otherwise}
  \end{cases}
\]
Contrary to what one might think, these functions $f_n$ never become identically
equal to $1$. Indeed, it can be observed that
\[
  \int_0^1 f_n(x)\, dx
  \le \sum_{k=1}^n \int_{q_k-r_k}^{q_k+r_k} 1\, dx
  = \sum_{k=1}^n \frac{1}{2\cdot 2^k}
  \le \frac{1}{2}\sum_{k=1}^{+\infty}\frac{1}{2^k} = \frac 1 2.
\]

It can be shown that $f_k$ is a Cauchy sequence in $H(0,1)$ and that it does not
converge in $H(0,1)$.
\end{example}

\begin{proof}
To verify that $f_n$ is a Cauchy sequence, we can simply observe that if $n\ge
m$, then $f_n$ and $f_m$ differ only on the intervals $I_k$ with $k$ between $m$
and $n$. Thus:
\[
  \int_0^1 \abs{f_n - f_m}
  \le \sum_{k=m+1}^{+\infty} \int_{q_k-r_k}^{q_k+r_k} 1\,dx
  = \sum_{k=m+1}^{+\infty}\frac{1}{2\cdot 2^k} = \frac{1}{2^{m+1}}.
\]
Since $f_n$ and $f_m$ take only the values $0$ and $1$, $\abs{f_n-f_m}$ also
takes only the values $0$ and $1$, so $\abs{f_n-f_m}^2=\abs{f_n-f_m}$. We have
thus verified that for $n\ge m$, we have:
\[
  \Abs{ f_n - f_m}_2 \le \sqrt{\frac{1}{2^{m+1}}}.
\]
It is therefore clear that for any chosen $\eps>0$, we can choose $N$ such that
$1/2^N \le \eps^2$, from which it follows that if $n,m\ge N$, then
$\Abs{f_n-f_m}_2 \le \eps$. That is: $f_n$ is a Cauchy sequence.

Now suppose, for contradiction, that there exists $f\in H(0,1)$ such that
$\Abs{f_n-f}_2 \to 0$. First of all, since each $f_k\le 1$, we can assume that
$f\le 1$ as well because otherwise, we could take $g(x) = \min\ENCLOSE{f(x),1}$
and would clearly have $\Abs{f_n - g}_2 \le \Abs{f_n -f}_2 \to 0$. In practice,
we are saying that by modifying the function $f$ without changing its integral,
we can assume that $f(x)\le 1$ for every $x\in[0,1]$.

Fixing an interval $I_n$, for $k\to +\infty$, we have:
\[
  0\le \int_{q_n-r_n}^{q_n+r_n}\abs{f_k(x)-f(x)}^2\, dx
  \le \int_0^1 \abs{f_k(x)-f(x)}^2\, dx = \Abs{ f_k -f}_2^2 \to 0.
\]
But now if $k\ge n$, it follows that $f_k=1$ on $I_n$, and therefore the value
of the integral above does not depend on $k$ and must therefore be identically
zero:
\[
  \int_{q_n-r_n}^{q_n+r_n} (1-f(x)) \, dx = 0
\]
and thus we can assert that $\sup f(I_n) = 1$ because if $\sup f(I_n) = \lambda
< 1$, the integral of $1-f(x)$ would be positive on the interval $I_n$.

We now want to demonstrate that on every interval $[a,b]\subset[0,1]$, if $b>a$,
we have $\sup f([a,b])=1$. Take $N$ large enough so that $r_N < (b-a)/3$. Then
the interval $[a+r_N, b-r_N]$ has a width greater than $r_N$ and contains
infinitely many rational numbers. There are therefore infinitely many indices
$n$ for which $q_n$ is in such an interval. Among these infinitely many, there
is certainly one with index $n\ge N$ (because the $q_n$ with $n<N$ are finite in
number). But if $q_n \in [a+r_N,b-r_N]$, then $I_n\subset[a,b]$, and therefore
$\sup f([a,b])\ge \sup f(I_n) = 1$.

This means that for every Riemann subdivision of the interval $[0,1]$, the
$\sup$ of $f$ on the subintervals of the subdivision is $1$, and therefore the
upper integral of $f$ is also $1$. Thus, being $f$ integrable,
\[
  \int_a^b f(x) = 1.
\]

We now want to conclude that this is in contradiction with the inequality
\[
  \int_a^b f_n(x) \le \frac 1 2
\]
that we observed at the beginning. Indeed, by the Cauchy-Schwarz inequality,
\begin{align*}
  \int_0^1 \abs{f(x)-f_n(x)}\, dx
  \le \Abs{f-f_n}_2\cdot \Abs{1}_2  \to 0
\end{align*}
and therefore, by the criterion of absolute convergence,
\[
  \int_0^1 (f(x) - f_n(x))\, dx \to 0.
\]
But we have seen that
\[
  \int_0^1 (f(x) - f_n(x))\, dx
  = \int_0^1 f(x)\, dx - \int_0^1 f_n(x)\, dx
  \ge 1 - \frac 1 2 = \frac 1 2
\]
thus obtaining a contradiction.
\end{proof}

The fact that we have found a Hilbertian basis for $H(0,2\pi)$ (and thus for
every $H(a,b)$) allows us to say that every function in $H(0,2\pi)$ can be
represented by its coordinates $c_k$ such that Bessel's equality holds:
\[
  \Abs{f}_2 = \sqrt{\sum_{k=0}^{+\infty} c_k^2}.
\]
We have indeed defined a correspondence:
\[
  \phi \colon H(0,2\pi) \to \ell^2,
  \qquad
  \ell^2=\ENCLOSE{c\in \RR^\NN \colon \sum_{k=0}^{+\infty} c_k^2<+\infty}
\]
defined by
\[
  c_k = \phi(f) = \langle f, e_k\rangle.
\]
Clearly, such $\phi$ is linear. Since $e_k$ is a Hilbertian basis, it is easy to
verify that $\phi$ is also injective. If on $\ell^2$, we define the inner
product and the corresponding norm:
\[
  \langle x,y\rangle = \sum_{k=0}^{+\infty} x_k y_k, \qquad
  \abs{x}_2 = \sqrt{\sum_{k=0}^{+\infty}} x_k^2
\]
Bessel's equality tells us that $\phi$ preserves the norm and thus the structure
of an Euclidean space. It could be shown that $\ell^2$ is complete, and this
means that $\phi$ is not surjective because the sequence $f_n$ considered in
example~\ref{ex:incompletezza_di_H} is Cauchy in $H(0,2\pi)$, thus $\phi(f_n)$
is Cauchy in $\ell^2$ and thus converges in $\ell^2$. The limit $c$ of the
sequence $\phi(f_n)$ represents a sequence of Fourier coefficients that cannot
be obtained from any function $f\in H(0,2\pi)$.

We can use Fourier series to solve the following\mynote{see historical notes on
page~\pageref{note:Basilea}}.

\begin{exercise}[Basel problem]
\label{ex:Basilea}%
We demonstrate that
\begin{equation}\label{eq:basilea}
  \sum_{k=1}^{+\infty} \frac{1}{k^2} = \frac{\pi^2}{6}.
\end{equation}
\end{exercise}
%%
%%
\begin{proof}[Solution.]
Let us calculate the Fourier coefficients of the function $f(x)=x$ in
$H(0,2\pi)$. We have
\[
  c_0 = \int_0^{2\pi} f(x)\cdot e_0(x)\, dx = \int_0^{2\pi}
 \frac{x}{\sqrt{2\pi}}\, dx =
 \frac{1}{\sqrt{2\pi}}\Enclose{\frac{x^2}{2}}_0^{2\pi}
  = \sqrt{2\pi^3}
\]
then
\begin{align*}
  c_{2k+1} &= \int_0^{2\pi} x \cdot \frac{\sin(kx)}{\sqrt \pi}\, dx \\
   &= \Enclose{\frac{x(-\cos(kx))}{k\sqrt \pi}}_0^{2\pi}
   - \int_0^{2\pi} \frac{-\cos(kx)}{k\sqrt \pi}\, dx\\
   &= \frac{-2\pi}{k\sqrt \pi} - 0 = -\frac{2\sqrt \pi}{k}
\end{align*}
and finally
\begin{align*}
  c_{2k} &= \int_0^{2\pi} x \cdot \frac{\cos (kx)}{\sqrt \pi}\, dx \\
  &= \Enclose{\frac{x \sin(kx)}{\sqrt \pi k}}_0^{2\pi}
  - \int_0^{2\pi} \frac{\sin(kx)}{\sqrt \pi k}\, dx
  = 0.
\end{align*}
Then Bessel's equation~\eqref{eq:uguaglianza_Bessel} tells us that
\[
  \Abs{x}_2^2 = \sum_{k=0}^{+\infty} c_k^2
  = c_0^2 + \sum_{k=0}^{+\infty} c_{2k+1}^2
  = 2\pi^3 + \sum_{k=1}^{+\infty} \frac{4\pi}{k^2}.
\]
But it is easy to calculate
\[
 \Abs{x}_2^2 = \int_0^{2\pi} x^2\, dx = \Enclose{\frac{x^3}{3}}_0^{2\pi}
  = \frac{8\pi^3}{3}
\]
and therefore
\[
  \frac{8\pi^3}{3} = 2\pi^3 + 4\pi \sum_{k=1}^{+\infty}\frac{1}{k^2}
\]
from which the desired result follows.
\end{proof}